
\documentclass[10pt,letterpaper]{article}
\usepackage[top=0.85in,left=2.75in,footskip=0.75in]{geometry}
\usepackage{amsmath,amssymb}
\usepackage{changepage}
\usepackage{textcomp,marvosym}
\usepackage{cite}
\usepackage{nameref,hyperref}
\usepackage[right]{lineno}
\usepackage[nopatch=eqnum]{microtype}
\DisableLigatures[f]{encoding = *, family = * }
\usepackage[table]{xcolor}
\usepackage{array}
\usepackage[aboveskip=1pt,labelfont=bf,labelsep=period,justification=raggedright,singlelinecheck=off]{caption}
\renewcommand{\figurename}{Fig}
\raggedright
\setlength{\parindent}{0.5cm}
\textwidth 5.25in
\textheight 8.75in
\linenumbers
\begin{document}
\begin{flushleft}
{\Large\bfseries PLOS ONE: Official PLOS Submission Preview}
\end{flushleft}
\section*{Abstract}
This submission-style preview is rendered from the official PLOS manuscript template package. The content remains placeholder text, but the page has been expanded into a more realistic article skeleton so the manuscript-style front matter, section organization, figure-caption handling, and end statements can be reviewed more professionally.
\section*{Author summary}
This placeholder summary demonstrates the front-matter flow expected in the PLOS manuscript template and makes the preview easier to compare with an actual submission manuscript.
\section{Introduction}
The verified preview uses the official PLOS manuscript-style setup. In accordance with the official template notes, figures are treated as separate submission assets rather than embedded graphics, so the page shows caption-based figure placement rather than inlined artwork.
\section{Related Work}
The related-work section is kept explicit so the preview resembles a real scientific manuscript rather than a single-page template sheet. This also helps reveal whether the section hierarchy remains clear when the article is read in manuscript form.
\section{Methodology}
The methodology section introduces a representative equation and a manuscript-style figure caption so the official PLOS package can be checked with more realistic technical material.
\begin{equation}
L = \sum_{i=1}^{M} w_i \left\| y_i - \hat{y}_i \right\|_2^2
\end{equation}
\paragraph{Fig 1.} Placeholder figure caption for a separately submitted figure asset showing the system pipeline, evaluation setup, or biological inspiration diagram.
\section{Experiments}
The experiments section provides enough structure to inspect the manuscript-style table placement and the body rhythm around quantitative content.
\subsection{Experimental Setup}
This paragraph acts as a stand-in for datasets, protocol details, baseline systems, and implementation notes.
\subsection{Quantitative Results}
\begin{table}[!ht]
\caption{Quantitative comparison block used to inspect the official PLOS manuscript template under a more realistic submission-style section flow.}
\centering
\begin{tabular}{lcc}
\hline
Method & Score & Time\\
\hline
Placeholder A & 0.91 & 1.4\\
Placeholder B & 0.88 & 1.7\\
Placeholder C & 0.84 & 2.0\\
\hline
\end{tabular}
\end{table}
\section{Conclusion and Discussion}
This page verifies the official PLOS manuscript template path for this journal target while preserving the manuscript-style conventions that matter most during initial submission review.
\section*{Data Availability}
Placeholder data availability statement.
\section*{References}
\begin{enumerate}
\item Placeholder A, Example B (2026) Verified template note for PLOS manuscript checking.
\end{enumerate}
\end{document}
